\chapter{2014年北京卷（理）}

\section{单选题}

\begin{problem}
已知集合 \(A = \{x \mid x^2 - 2x = 0\}\)，\(B = \{0, 1, 2\}\)，则 \(A \cap B =\)（\quad）

\choices
    {\( \{0\} \)}
    {\(\{0, 1\}\)}
    {\(\{0, 2\}\)}
    {\(\{0, 1, 2\}\)}
\end{problem}

\begin{problem}
下列函数中，在区间 \((0, +\infty)\) 上为增函数的是（\quad）

\choices
    {\(y = \sqrt{x + 1}\)}
    {\( y = (x - 1)^{2} \)}
    {\( y = 2^{-x} \)}
    {\( y = \log_{0.5}(x + 1) \)}
\end{problem}

\begin{problem}
曲线 \(\left\{ \begin{array}{l} x = -1 + \cos \theta \\ y = 2 + \sin \theta. \end{array} \right.\) （\(\theta\) 为参数）的对称中心（\quad）

\choices
    {在直线 \( y = 2x \) 上}
    {在直线 \( y = -2x \) 上}
    {在直线 \( y = x - 1 \) 上}
    {在直线 \( y = x + 1 \) 上}
\end{problem}

\begin{problem}
当 \(m = 7\)，\(n = 3\) 时，执行如图所示的程序框图，输出的 \(S\) 值为（\quad）

\texfigure[width=0.42\linewidth,max width=0.42\linewidth]{img_repaint/2014/beijing_science/q4_fig1.tex}

\choices
    {7}
    {42}
    {210}
    {840}
\end{problem}

\begin{problem}
设 \(\{a_{n}\}\) 是公比为 \(q\) 的等比数列，则 \(q > 1\) 是“\(\{a_{n}\}\) 为递增数列”的（\quad）

\choices
    {充分且不必要条件}
    {必要且不充分条件}
    {充分必要条件}
    {既不充分也不必要条件}
\end{problem}

\begin{problem}
若 \(x\)，\(y\) 满足 \(\begin{cases}
x + y - 2 \geqslant 0,\\
kx - y + 2 \geqslant 0,\\
y \geqslant 0,
\end{cases}\) 且 \(z = y - x\) 的最小值为 \(-4\)，则 \(k\) 的值为（\quad）

\choices
    {2}
    {\(-2\)}
    {\( \frac{1}{2} \)}
    {\(-\frac{1}{2}\)}
\end{problem}

\begin{problem}
在空间直角坐标系 \(Oxyz\) 中，已知 \(A(2,0,0)\)，\(B(2,2,0)\)，\(C(0,2,0)\)，\(D(1,1,\sqrt{2})\)，若 \(S_{1}\)，\(S_{2}\)，\(S_{3}\) 分别表示三棱锥 \(D-ABC\) 在 \(xOy\)，\(yOz\)，\(zOx\) 坐标平面上的正投影图形的面积，则（\quad）

\choices
    {\(S_{1} = S_{2} = S_{3}\)}
    {\(S_{1} = S_{2}\) 且 \(S_{3}\neq S_{1}\)}
    {\(S_{1} = S_{3}\) 且 \(S_{3}\neq S_{2}\)}
    {\(S_{2} = S_{3}\) 且 \(S_{1}\neq S_{3}\)}
\end{problem}

\begin{problem}
学生的语文，数学成绩均被评定为三个等级，依次为“优秀”“合格”“不合格”. 若学生甲的语文，数学成绩都不低于学生乙，且其中至少有一门成绩高于乙，则称“学生甲比学生乙成绩好”. 如果一组学生中没有哪位学生比另一位学生成绩好，并且不存在语文成绩相同，数学成绩也相同的两位学生，则这一组学生最多有（\quad）

\choices
    {2 人}
    {3 人}
    {4 人}
    {5 人}
\end{problem}

\section{填空题}

\begin{problem}
复数 \(\left(\frac{1 + \i}{1 - \i}\right)^2 =\)\fillinblank{}.
\end{problem}

\begin{problem}
已知向量 \(a\)，\(b\) 满足 \(|a| = 1\)，\(b = (2, 1)\)，且 \(\lambda a + b = 0\)（\(\lambda \in \R\)），则 \(|\lambda| = \)\fillinblank{}.
\end{problem}

\begin{problem}
设双曲线 \(C\) 经过点 \((2,2)\)，且与 \(\frac{y^{2}}{4}-x^{2}=1\) 具有相同渐近线，则 \(C\) 的方程为\fillinblank{}；渐近线方程为\fillinblank{}.
\end{problem}

\begin{problem}
若等差数列 \(\{a_{n}\}\) 满足 \(a_{7} + a_{8} + a_{9} > 0\)，\(a_{7} + a_{10} < 0\)，则当 \(n =\fillinblank{}\) 时，\(\{a_{n}\}\) 的前 \(n\) 项和最大.
\end{problem}

\begin{problem}
把 5 件不同产品摆成一排，若产品 \(A\) 与产品 \(B\) 相邻，且产品 \(A\) 与产品 \(C\) 不相邻，则不同的摆法有\fillinblank{} 种.
\end{problem}

\begin{problem}
设函数 \( f(x) = A \sin (\omega x + \varphi) \) （\(A\)，\(\omega\)，\(\varphi\) 是常数，\( A > 0 \)，\(\omega > 0\)）.若 \( f(x) \) 在区间 \( \left[\frac{\pi}{6}, \frac{\pi}{2}\right] \) 上具有单调性，且 \( f\left(\frac{\pi}{2}\right) = f\left(\frac{2\pi}{3}\right) = -f\left(\frac{\pi}{6}\right) \)，则 \( f(x) \) 的最小正周期为\fillinblank{}.
\end{problem}

\section{解答题}

\begin{problem}
如图，在 \(\triangle ABC\) 中，\(\angle B = \frac{\pi}{3}\)，\(AB = 8\)，点 \(D\) 在 \(BC\) 上，且 \(CD = 2\)，\(\cos \angle ADC = \frac{1}{7}\).

\begin{enumerate}
\item 求 \(\sin \angle BAD\).
\item 求 \(BD\)，\(AC\) 的长.
\end{enumerate}

\bitmapfigure[max width=0.38\linewidth]{img/2014/beijing_science/q15_fig1.png}

\end{problem}

\begin{problem}
李明在 10 场篮球比赛中的投篮情况（假设各场比赛相互独立）：

\begin{center}
\begin{tabular}{c|c|c|c|c|c}
\hline
场次 & 投篮次数 & 命中次数 & 场次 & 投篮次数 & 命中次数 \\
\hline
主场 1 & 22 & 12 & 客场 1 & 18 & 8 \\
主场 2 & 15 & 12 & 客场 2 & 13 & 12 \\
主场 3 & 12 & 8 & 客场 3 & 21 & 7 \\
主场 4 & 23 & 8 & 客场 4 & 18 & 15 \\
主场 5 & 24 & 20 & 客场 5 & 25 & 12 \\
\hline
\end{tabular}
\end{center}

\begin{enumerate}
\item 从上述比赛随机选择一场，求李明在该场比赛中的投篮命中率超过 0.6 的概率；
\item 从上述比赛中随机选择一个主场和客场，求李明的投篮命中率一场超过 0.6，一场不超过 0.6 的概率；
\item 记 \(x\) 是表中 10 个命中次数的平均数，从上述比赛中随机选择一场，记 \(X\) 为李明在这场比赛中命中次数，比较 \(E(X)\) 与 \(x\) 的大小. （只需要写出结论）
\end{enumerate}
\end{problem}

\begin{problem}
如图，正方形 \(AMDE\) 的边长为 2，\(B\)，\(C\) 分别为 \(AM\)，\(MD\) 的中点，在五棱锥 \(P\)-\(ABCDE\) 中，\(F\) 为棱 \(PE\) 的中点，平面 \(ABF\) 与棱 \(PD\)，\(PC\) 分别交于点 \(G\)，\(H\).

\begin{enumerate}
\item 求证：\(AB \parallel FG\).
\item 若 \(PA \perp\) 平面 \(ABCDE\)，且 \(PA = AE\)，求直线 \(BC\) 与平面 \(ABF\) 所成角的大小，并求线段 \(PH\) 的长.
\end{enumerate}

\bitmapfigure[max width=0.42\linewidth]{img/2014/beijing_science/q17_fig1.png}

\end{problem}

\begin{problem}
已知函数 \(f(x) = x\cos x - \sin x\)，\(x \in \left[0, \frac{\pi}{2}\right]\).

\begin{enumerate}
\item 求证：\(f(x) \leqslant 0\).
\item 若 \(a < \frac{\sin x}{x} < b\) 在 \(\left(0, \frac{\pi}{2}\right)\) 上恒成立，求 \(a\) 的最大值与 \(b\) 的最小值.
\end{enumerate}
\end{problem}

\begin{problem}
已知椭圆 \(C: x^{2} + 2y^{2} = 4\).

\begin{enumerate}
\item 求椭圆 \(C\) 的离心率.
\item 设 \(O\) 为坐标原点，若点 \(A\) 在椭圆 \(C\) 上，点 \(B\) 在直线 \(y = 2\) 上，且 \(OA \perp OB\)，试判断直线 \(AB\) 与圆 \(x^{2} + y^{2} = 2\) 的位置关系，并证明你的结论.
\end{enumerate}
\end{problem}

\begin{problem}
对于数对序列 \(P: (a_1, b_1)\)，\((a_2, b_2)\)，\(\cdots\)，\((a_n, b_n)\)，记 \(T_1(P) = a_1 + b_1\)，\(T_k(P) = b_k + \max \{T_{k-1}(P), a_1 + a_2 + \cdots + a_k\} (2 \leqslant k \leqslant n)\)，其中 \(\max \{T_{k-1}(P), a_1 + a_2 + \cdots + a_k\}\) 表示 \(T_{k-1}(P)\) 和 \(a_1 + a_2 + \cdots + a_k\) 两个数中最大的数.

\begin{enumerate}
\item 对于数对序列 \(P: (2,5)\)，\((4,1)\)，求 \(T_{1}(P)\)，\(T_{2}(P)\) 的值.
\item 记 \(m\) 为 \(a\)，\(b\)，\(c\)，\(d\) 四个数中最小值，对于由两个数对 \((a, b)\)，\((c, d)\) 组成的数对序列 \(P: (a, b)\)，\((c, d)\) 和 \(P': (c, d)\)，\((a, b)\).试分别对 \(m = a\) 和 \(m = d\) 时两种情况比较 \(T_{2}(P)\) 和 \(T_{2}(P')\) 的大小.
\item 在由 5 个数对 \((11,8)\)，\((5,2)\)，\((16,11)\)，\((11,11)\)，\((4,6)\) 组成的所有数对序列中，写出一个数对序列 \(P\) 使 \(T_{5}(P)\) 最小，并写出 \(T_{5}(P)\) 的值. （只需写出结论）
\end{enumerate}
\end{problem}
